\documentclass[a4paper,12pt]{article}
\usepackage[utf8]{inputenc}
\usepackage[T1]{fontenc}
\usepackage[ngerman]{babel}
\usepackage{amsmath}
\usepackage{amssymb}
\usepackage{enumitem}
\usepackage{geometry}
\geometry{a4paper, margin=2.5cm}
\usepackage{parskip}

\title{Einsendeaufgaben -- Aufgabe 5.1\\[0.5em]
\large 61211 Analysis}
\author{}
\date{}

\begin{document}
\maketitle

\section*{Aufgabe 1}

\begin{enumerate}[label=\alph*), start=2]

    \item
    Wir nutzen das Majorantenkriterium.

    Wir definieren
    \[
        f: (1,\infty) \to \mathbb{R} : (x+\sin^2(x))e^{-x}
        \quad\text{und}\quad
        g: (1,\infty) \to \mathbb{R} : (x+1)e^{-x}
    \]

    Zunächst zeigen wir, dass $\int_1^\infty g(x)\,dx = \int_1^\infty (x+1)e^{-x}\,dx$ integrierbar ist. Für ein $a \in [1,\infty)$ gilt:
    \begin{align*}
        \int_1^a (x+1)e^{-x}\,dx
        &= (x+1)(-e^{-x})\Big|_1^a - \int_1^a e^{-x} \\
        &= -\frac{(a-1)}{e^a} + 2e^{-1} + e^{-a} - e^{-1} \\
        &= -\frac{a+1}{e^a} + e^{-a} + e^{-1}
    \end{align*}

    Es gilt $\lim_{a\to\infty} e^{-a} = 0$. Außerdem existiert der Grenzwert von $\lim_{a\to\infty} \frac{a+1}{e^a}$, da $0 \leq \frac{a+1}{e^a} \leq 1$ wegen $e^a \geq a+1$.

    Demnach ist $\int_1^\infty g(x)\,dx$ integrierbar.

    Für alle $x \in (1,\infty)$ gilt
    \begin{align*}
        |f(x)| &= \left| (x+\sin^2(x))e^{-x} \right| \\
        &= (x+\sin^2(x))e^{-x} \qquad (\text{wegen } x>1 \text{ und } 0 \leq \sin^2(x) \leq 1) \\
        &\leq (x+1)e^{-x} \\
        &= g(x).
    \end{align*}

    Da $f$ stetig ist, gilt auch $f \in R([\alpha,\beta])$ für alle $1 < \alpha < \beta < \infty$. Wir können also das Majorantenkriterium anwenden. Demnach ist $\int_1^\infty (x+\sin^2(x))e^{-x}\,dx$ integrierbar.

\end{enumerate}

\end{document}
